\chapter{جمع‌بندی پروژه بازنگری مقاله متاسطح}

\section{هدف گزارش}
هدف این گزارش، ایجاد یک تصویر روشن و منسجم از وضعیت فعلی مقاله، روش مدل‌سازی، کدهای موجود، نقدهای داوران و مسیر پیشنهادی برای آماده‌سازی مقاله جهت ارسال به یک مجله معتبر در حوزه نانوفوتونیک یا فیزیک کاربردی است. مقاله فعلی درباره پیش‌بینی پاسخ نوری متاسطح‌ها، شامل دامنه و فاز نور عبوری، با استفاده از یک شبکه عصبی مقید به فیزیک است. ایده اصلی مقاله ارزشمند است، اما نسخه فعلی برای پذیرش در ژورنال معتبر نیاز به بازنگری علمی، تجربی و نگارشی دارد.



\section{خلاصه ایده علمی مقاله}
مقاله یک شبکه عصبی برای پیش‌بینی پاسخ نوری واحدهای سازنده متاسطح پیشنهاد می‌کند. ورودی مدل شامل پارامترهای هندسی و فیزیکی ساختار است؛ از جمله قطر نانودیسک، ارتفاع، دوره تناوب، طول موج، قطبش، زاویه تابش، ضریب شکست ماده لنز و ضریب شکست زیرلایه. خروجی مدل دامنه عبور و فاز نور عبوری است.

چالش اصلی در این مسئله، رفتار پیچیده فاز است. فاز یک کمیت دوری است و در بازه \lr{$0$} تا \lr{$2\pi$} تعریف می‌شود. در نتیجه دو مقدار عددی که ظاهرا از هم دورند، مثل \lr{$0$} و \lr{$2\pi$}، از نظر فیزیکی می‌توانند بسیار نزدیک یا حتی معادل باشند. اگر فاز مستقیما به صورت یک عدد اسکالر پیش‌بینی شود، شبکه ممکن است در نزدیکی مرزهای دوری دچار خطای بزرگ ظاهری شود. برای مثال، اگر مقدار واقعی فاز نزدیک صفر باشد و مدل مقداری نزدیک \lr{$2\pi$} پیش‌بینی کند، خطای عددی مستقیم بزرگ به نظر می‌رسد، در حالی که خطای فیزیکی کوچک است.

\section{داده‌ها و مسئله فیزیکی}
داده‌ها از شبیه‌سازی‌های \lr{FEM} با نرم‌افزار \lr{JCMsuite} به دست آمده‌اند. ساختار مورد مطالعه یک آرایه دوبعدی همگن از نانودیسک‌هاست که هر سلول واحد شامل یک نانودیسک است. با اعمال شرایط مرزی متقارن، پاسخ یک آرایه بی‌نهایت از نانودیسک‌ها شبیه‌سازی شده است. مقاله سه اندازه داده را بررسی می‌کند:

\begin{itemize}
	\item مجموعه داده \lr{A}: شامل \lr{1075} نمونه؛
	\item مجموعه داده \lr{B}: شامل \lr{580} نمونه؛
	\item مجموعه داده \lr{C}: شامل \lr{290} نمونه و نماینده حالت کمبود شدید داده.
\end{itemize}

این سه مجموعه داده برای نشان دادن رفتار مدل در مقیاس‌های مختلف داده استفاده شده‌اند. در نتایج گزارش‌شده، مقدار \lr{$R^2$} برای فاز در مجموعه \lr{A} برابر حدود \lr{0.928}، در مجموعه \lr{B} برابر حدود \lr{0.906} و در مجموعه \lr{C} برابر حدود \lr{0.859} است. خطای زاویه‌ای میانگین نیز برای مجموعه‌های \lr{A} و \lr{B} به ترتیب حدود \lr{0.114} و \lr{0.124} رادیان گزارش شده است.



\section{معماری مدل در کد}
فایل اصلی معماری مدل در مسیر \lr{src/models/base\_model.py} قرار دارد. مدل یک شبکه عصبی \lr{TensorFlow/Keras} است که ورودی \lr{8} بعدی دریافت می‌کند. ابتدا ورودی‌ها نرمال‌سازی می‌شوند. قطر نسبت به دوره تناوب نرمال می‌شود تا اطلاعات مربوط به نسبت پرشدگی یا اندازه نسبی نانودیسک نسبت به سلول واحد حفظ شود. ارتفاع، دوره تناوب و زاویه تابش نیز با قواعد ساده نرمال‌سازی می‌شوند. طول موج و ضرایب شکست در داده فعلی ثابت‌اند و در صورت نبود تغییر، مقدار ثابت برای آن‌ها در نظر گرفته شده است.

بدنه شبکه از چند لایه متراکم تشکیل شده است. ابتدا یک لایه با \lr{256} نورون، سپس لایه‌های \lr{128} و \lr{64} نورونی استفاده شده‌اند. بعد از لایه‌های متراکم، \lr{Batch Normalization}، فعال‌ساز \lr{ReLU} و در برخی بخش‌ها \lr{Dropout} قرار گرفته است. در کد، ایده اتصال میان‌بر نیز توضیح داده شده اما در نسخه فعلی غیرفعال است.

مدل دو شاخه خروجی دارد:


این بخش، ستون فقرات نوآوری مقاله است و باید در متن مقاله بسیار شفاف‌تر و دقیق‌تر توضیح داده شود. بهتر است گفته شود این لایه پارامتر trainable ندارد، اما عملیات آن در مسیر مشتق‌گیری و آموزش حضور دارد. عبارت \lr{non-trainable layer} به تنهایی می‌تواند سوءبرداشت ایجاد کند، چون برخی داوران آن را با فعال‌سازها یا سایر اجزای قابل یادگیری شبکه مقایسه می‌کنند.

\section{تابع زیان و آموزش}
تابع زیان در فایل \lr{src/losses/basic\_losses.py} تعریف شده است. سه مولفه اصلی دارد:

\begin{itemize}
	\item زیان دامنه: میانگین مربعات اختلاف دامنه واقعی و پیش‌بینی‌شده؛
	\item زیان فاز: فاصله اقلیدسی بین نقطه واقعی و نقطه پیش‌بینی‌شده روی دایره واحد؛
	\item زیان انرژی: جریمه برای حالتی که دامنه پیش‌بینی‌شده از \lr{1} بزرگ‌تر شود.
\end{itemize}

زیان فاز از نظر مفهومی مناسب است، چون به جای مقایسه مستقیم دو عدد فاز، فاصله دو بردار روی دایره واحد را اندازه می‌گیرد:
\[
L_{\mathrm{phase}}=\frac{1}{N}\sum_i \sqrt{(\cos\phi_i-\hat{c}_i)^2+(\sin\phi_i-\hat{s}_i)^2}
\]
این معیار با ماهیت دوری فاز سازگار است و مشکل پرش مصنوعی در مرز \lr{$0/2\pi$} را کاهش می‌دهد.

اما زیان انرژی در نسخه فعلی یک نقطه ضعف مهم است. شاخه دامنه در معماری مدل از فعال‌ساز \lr{sigmoid} استفاده می‌کند؛ بنابراین خروجی دامنه ذاتا در بازه \lr{$[0,1]$} قرار می‌گیرد. از طرف دیگر زیان انرژی فقط زمانی مقدار غیرصفر می‌گیرد که دامنه پیش‌بینی‌شده بزرگ‌تر از \lr{1} باشد:


\section{وزن‌دهی بر اساس شکاف}
در تابع زیان، وزن نمونه‌ها بر اساس اندازه شکاف محاسبه می‌شود. اگر شکاف بزرگ‌تر از آستانه باشد، وزن \lr{1} می‌گیرد و اگر کوچک‌تر باشد، وزن \lr{small\_gap\_weight} می‌گیرد. در توضیحات اولیه کد، این ایده به این صورت مطرح شده که نمونه‌های شکاف کوچک، به دلیل رفتار پیچیده‌تر، در طول آموزش وزن متفاوت داشته باشند.

با این حال، وضعیت فعلی این بخش نیاز به بازبینی دارد. در فایل \lr{trainer.py} مقدارهای مربوط به وزن شکاف کوچک طوری تنظیم شده‌اند که عملا تغییر معنی‌داری در طول آموزش رخ نمی‌دهد؛ مثلا شروع \lr{1.0} و پایان \lr{0.9} با نقطه شروع ramp در \lr{1000} epoch، در حالی که آموزش اصلی \lr{400} epoch است. بنابراین در نسخه فعلی، وزن‌دهی پویا بیشتر در حد ایده و توضیح باقی مانده است تا یک مولفه فعال در آزمایش.

برای مقاله، دو راه وجود دارد. یا این بخش باید به صورت جدی و کنترل‌شده فعال و گزارش شود، یا باید از متن اصلی حذف/کمرنگ شود تا داور احساس نکند یک مولفه ادعایی بدون نقش واقعی معرفی شده است. اگر قرار باشد استفاده شود، باید آزمایش ablation داشته باشد: مدل بدون وزن‌دهی شکاف، مدل با وزن‌دهی ثابت، و مدل با زمان‌بندی وزن‌دهی.

\section{نتایج فعلی و تفسیر آن‌ها}
نتایج فعلی نشان می‌دهد مدل روی مجموعه‌های \lr{A} و \lr{B} عملکرد خوبی دارد. خطای زاویه‌ای حدود \lr{0.11} تا \lr{0.12} رادیان برای مسئله طراحی متاسطح قابل توجه است، به‌خصوص اگر با هزینه بالای تولید داده‌های شبیه‌سازی مقایسه شود. همچنین افت عملکرد در مجموعه \lr{C} طبیعی است و می‌تواند به عنوان نشانه‌ای از محدودیت داده بسیار کم تفسیر شود.


\section{نقدهای داوران و معنای آن‌ها}
مقاله دو بار رد شده است. رد شدن اول در \lr{Scientific Reports} و رد شدن دوم در \lr{Machine Learning: Science and Technology} بوده است. نقدهای داوران را نباید فقط به عنوان پاسخ به همان ژورنال‌ها دید، چون هدف فعلی ارسال مجدد به همان داوران نیست. بهتر است این نقدها به عنوان نشانه‌هایی از نقاط آسیب‌پذیر مقاله در برابر هر داور دقیق استفاده شوند.

نقدهای اصلی را می‌توان در چند دسته خلاصه کرد.

\subsection{ادعاهای زائد یا گمراه‌کننده}
مهم‌ترین مورد، انرژی‌لاس است. وقتی خروجی دامنه با \lr{sigmoid} محدود می‌شود، جریمه دامنه بزرگ‌تر از یک عملا نقشی در آموزش ندارد. بنابراین اگر مقاله ادعا کند مدل به دلیل energy penalty قانون بقای انرژی را یاد می‌گیرد، داور می‌تواند این ادعا را رد کند. بهتر است گفته شود خروجی دامنه به صورت معماری در بازه فیزیکی مجاز محدود شده است. اگر energy penalty باقی می‌ماند، باید نقش آن فقط در آزمایش بدون \lr{sigmoid} یا در نسخه جایگزین مدل توضیح داده شود.

مورد دوم، اصطلاح \lr{non-trainable} است. نرمال‌سازی فاز پارامتر قابل یادگیری ندارد، اما در مسیر آموزش و backpropagation حضور دارد. بنابراین باید با دقت گفته شود: «یک نگاشت قطعی و بدون پارامتر یادگیری‌پذیر که خروجی خام فاز را روی دایره واحد تصویر می‌کند.»

مورد سوم، ارجاع و توصیف کار \lr{An et al.} است. در نسخه فعلی گفته شده آن کار از دو شبکه جداگانه استفاده کرده، در حالی که طبق نقد داور، آن مقاله یک شبکه با دو head داشته است. این خطا باید حتما اصلاح شود، چون خطای ارجاعی اعتماد داور به دقت نویسندگان را کاهش می‌دهد.

\subsection{نبود مقایسه کافی}
بزرگ‌ترین ضعف مقاله نبود baselineهای مناسب است. نسخه فعلی بیشتر نشان می‌دهد مدل پیشنهادی خوب کار می‌کند، اما نشان نمی‌دهد چرا این مدل از گزینه‌های ساده‌تر بهتر یا مناسب‌تر است. حداقل baselineهای لازم عبارت‌اند از:


این مقایسه‌ها نباید بی‌هدف و شبیه benchmark عمومی یادگیری ماشین طراحی شوند. هدف آن‌ها باید پاسخ به پرسش‌های فیزیکی و روشی مقاله باشد: آیا نمایش دوری فاز لازم است؟ آیا نرمال‌سازی روی دایره واحد واقعا خطای فاز را کم و خروجی‌ها را فیزیکی‌تر می‌کند؟ آیا مدل پیشنهادی نسبت به نمایش حقیقی/موهومی مزیت دارد؟

\subsection{طراحی آماری و داده‌های کوچک}
تقسیم تصادفی ساده در داده‌های کوچک و نامتوازن کافی نیست. به‌خصوص برای مجموعه \lr{C}، اگر فقط \lr{10\%} داده برای آزمون استفاده شود، تعداد نمونه آزمون حدود \lr{29} می‌شود. چنین آزمونی برای گزارش \lr{$R^2$} و نتیجه‌گیری قوی بسیار ناپایدار است. در نسخه اصلاح‌شده باید از \lr{k-fold cross-validation}، به‌ویژه برای مجموعه \lr{C}، استفاده شود. همچنین بهتر است تقسیم داده‌ها بر اساس بازه‌های شکاف و شاید فاز طبقه‌بندی شود تا هر fold نماینده مناسبی از نواحی فیزیکی مهم باشد.

\subsection{جزئیات فیزیکی ناقص}
داور اول \lr{Scientific Reports} اشاره کرده بود که جزئیات طراحی متاسطح کافی نیست؛ از جمله ماده، طول موج، هندسه سلول واحد و شرایط شبیه‌سازی. در نسخه فعلی درفت، برخی از این موارد به Supplementary Material ارجاع داده شده‌اند، اما برای مقاله اصلی کافی نیست. خواننده باید در همان متن اصلی بفهمد ساختار چیست، ماده چیست، طول موج کاری چیست، بازه قطر/ارتفاع/دوره تناوب چیست و شرایط مرزی یا نوع شبیه‌سازی چگونه بوده است.

\section{جایگاه‌یابی مقاله}
یکی از تصمیم‌های راهبردی مهم این است که مقاله نباید به عنوان یک مقاله اصلی در حوزه یادگیری ماشین معرفی شود. مقاله یک روش جدید بنیادی در \lr{ML} ارائه نمی‌کند؛ بلکه از یک طراحی شبکه عصبی با قیدهای فیزیکی/هندسی برای حل یک مسئله نانوفوتونیکی استفاده می‌کند. بنابراین چارچوب درست مقاله باید چنین باشد:

\begin{quote}
	یک مدل داده‌کارآمد و فیزیکی‌معنادار برای پیش‌بینی پاسخ نوری متاسطح‌ها در شرایط کمبود داده شبیه‌سازی.
\end{quote}

این framing باید از عنوان، چکیده و پاراگراف‌های اول مقدمه مشخص باشد. اصطلاح \lr{physics-informed neural network} در نسخه فعلی ممکن است انتظاراتی ایجاد کند که مقاله برآورده نمی‌کند، چون در برخی جوامع علمی، \lr{PINN} به مدل‌هایی گفته می‌شود که معادلات دیفرانسیل جزئی مانند معادلات ماکسول را مستقیما در تابع زیان وارد می‌کنند. مدل حاضر بیشتر \lr{physics-constrained} یا \lr{physics-guided} است، چون قیدهای فیزیکی/هندسی خروجی را در معماری و تابع زیان لحاظ می‌کند.

بنابراین پیشنهاد می‌شود از واژه \lr{physics-constrained} به عنوان اصطلاح اصلی استفاده شود و اگر \lr{physics-informed} به کار می‌رود، بسیار دقیق تعریف شود. هدف این نیست که نقد داور سوم را بپذیریم که «این کار واقعا physics-informed نیست»، بلکه باید از ابتدا تعریف درست خودمان را روشن کنیم تا داور در چارچوب اشتباه قضاوت نکند.

\section{ضعف‌های فعلی نسخه درفت}
نسخه فعلی درفت از نظر علمی پایه خوبی دارد، اما چند ضعف جدی دارد:



\section{نقشه راه پیشنهادی برای ادامه پروژه}
برای آماده‌سازی مقاله جهت ارسال مجدد، بهتر است کار به چند فاز تقسیم شود.

\subsection{فاز اول: تصمیم‌های راهبردی متن}
در این فاز باید تصمیم گرفته شود مقاله با چه ادعای مرکزی بازنویسی شود. پیشنهاد این است که ادعای مرکزی چنین باشد:

\begin{quote}
	مدل پیشنهادی با استفاده از نمایش هندسی فاز روی دایره واحد و محدودسازی فیزیکی دامنه، امکان پیش‌بینی داده‌کارآمد پاسخ نوری متاسطح را با مجموعه داده‌های کوچک و نامتوازن فراهم می‌کند.
\end{quote}

در این چارچوب، انرژی‌لاس نباید به عنوان ستون اصلی نوآوری مطرح شود. ستون اصلی، نمایش فاز و نرمال‌سازی دایره واحد است. دامنه نیز با \lr{sigmoid} در بازه فیزیکی محدود می‌شود.

\subsection{فاز دوم: اصلاح و تکمیل آزمایش‌ها}
در این فاز باید کدها برای اجرای baselineها و ارزیابی آماری بهتر توسعه داده شوند. آزمایش‌های پیشنهادی عبارت‌اند از:

\begin{enumerate}
	\item مدل پیشنهادی فعلی با خروجی دامنه و سینوس/کسینوس نرمال‌شده؛
	\item مدل بدون نرمال‌سازی فاز؛
	\item مدل با پیش‌بینی مستقیم فاز؛
	\item مدل sine-only؛
	\item مدل real/imaginary برای ضریب عبور؛
	\item ارزیابی همه مدل‌ها با \lr{k-fold cross-validation} و تقسیم طبقه‌بندی‌شده.
\end{enumerate}

برای هر مدل باید معیارهای یکسان گزارش شوند: خطای دامنه، خطای فاز در فضای دایره‌ای، خطای زاویه‌ای، \lr{$R^2$} مناسب برای فاز، و در صورت امکان عملکرد تفکیک‌شده در نواحی شکاف کوچک و بزرگ.


\section{جمع‌بندی نهایی}
این پروژه از نظر ایده اصلی ارزشمند است. مسئله پیش‌بینی فاز و دامنه در متاسطح‌ها مسئله‌ای واقعی و مهم است، و استفاده از نمایش سینوس/کسینوس به همراه نرمال‌سازی روی دایره واحد، راه‌حلی طبیعی و قابل دفاع برای مشکل دوری بودن فاز است. نتایج فعلی نیز نشان می‌دهد مدل می‌تواند با داده کم عملکرد خوبی داشته باشد.

اما برای رسیدن به سطح قابل انتشار، مقاله باید از حالت «گزارش یک مدل موفق» به حالت «استدلال علمی کامل و قابل دفاع» تبدیل شود. این یعنی ادعاهای زائد حذف یا اصلاح شوند، baselineهای لازم اضافه شوند، طراحی ارزیابی آماری تقویت شود، و متن مقاله از ابتدا در چارچوب نانوفوتونیک و فیزیک کاربردی نوشته شود.

اولویت فوری پروژه، اجرای مقایسه‌های پایه و حل مسئله energy loss است. تا وقتی این دو مورد حل نشوند، بازنویسی کامل متن ممکن است زودهنگام باشد. پس از آن، می‌توان با اطمینان بیشتری چکیده، مقدمه، روش و نتایج را بازنویسی کرد و مقاله را برای یک مجله مناسب‌تر در حوزه نانوفوتونیک آماده ساخت.